\documentclass[12pt,a4paper]{article}

\usepackage[utf8]{inputenc}
\usepackage[T1]{fontenc}
\usepackage[brazil]{babel}
\usepackage{geometry}
\usepackage{graphicx}
\usepackage{booktabs}
\usepackage{longtable}
\usepackage{array}
\usepackage{float}
\usepackage{hyperref}
\usepackage{url}
\usepackage{caption}
\usepackage{enumitem}

\geometry{margin=2.5cm}
\graphicspath{{figuras/}}
\setlength{\parskip}{0.7em}
\setlength{\parindent}{0pt}

\title{GraphQL vs REST: Um Experimento Controlado}
\author{Laboratório 05 -- Experimentação de Software}
\date{25 de junho de 2026}

\begin{document}

\maketitle

\begin{abstract}
Este relatório apresenta um experimento controlado para comparar APIs REST e GraphQL em duas métricas quantitativas: tempo de resposta, medido em milissegundos, e tamanho da resposta, medido em bytes. O experimento foi executado em ambiente Docker para reduzir variações locais de runtime e dependências. Foram coletadas 1000 medições para REST e 1000 medições para GraphQL, totalizando 2000 registros válidos e 0 falhas. Os resultados indicam que GraphQL reduziu substancialmente o tamanho do payload, mas não foi mais rápido que REST no cenário avaliado.
\end{abstract}

\section{Versão resumida}

\textbf{Introdução.} O estudo compara REST e GraphQL em uma API local com base fictícia de estatísticas de futebol.

\textbf{Problema.} Não está claro se GraphQL, por permitir selecionar campos específicos, entrega respostas mais rápidas e menores que REST em um cenário controlado.

\textbf{Questões de pesquisa.} RQ1: respostas GraphQL são mais rápidas que respostas REST? RQ2: respostas GraphQL têm tamanho menor que respostas REST?

\textbf{Hipóteses.} Para RQ1, a hipótese alternativa afirma que GraphQL tem menor tempo médio que REST. Para RQ2, a hipótese alternativa afirma que GraphQL tem menor tamanho médio de resposta que REST.

\textbf{Metodologia.} Foram feitas requisições pareadas para os mesmos IDs de jogadores, alternando os tratamentos REST e GraphQL. A coleta oficial teve 1000 iterações por tratamento.

\textbf{Resultado principal.} REST apresentou mediana de tempo menor, portanto RQ1 não foi confirmada. GraphQL apresentou payload muito menor, confirmando RQ2.

\textbf{Conclusão.} Neste experimento, GraphQL foi vantajoso para reduzir tráfego de dados, mas não para reduzir tempo de resposta.

\section{Introdução}

APIs Web são elementos centrais na integração entre aplicações. Entre as abordagens mais utilizadas, REST organiza a comunicação por recursos e endpoints, seguindo princípios arquiteturais consolidados por Fielding \cite{fielding2000rest}. GraphQL, por sua vez, define uma linguagem de consulta baseada em schemas, em que o cliente informa explicitamente os campos desejados na resposta \cite{graphqlspec2025}.

Segundo a documentação da GraphQL Foundation, uma consulta GraphQL retorna uma estrutura com o mesmo formato da consulta e permite que o servidor saiba exatamente quais campos o cliente está solicitando \cite{graphqlqueries2026}. Essa característica sugere uma possível redução do tamanho das respostas, especialmente quando o cliente precisa de apenas parte dos dados disponíveis em um recurso.

Este relatório documenta o Laboratório 05, cujo objetivo foi avaliar quantitativamente diferenças entre REST e GraphQL em um experimento controlado.

\section{Contextualização}

O sistema experimental implementa duas interfaces sobre a mesma base de dados fictícia de estatísticas de futebol:

\begin{itemize}
  \item uma API REST, com endpoint para consultar jogadores;
  \item uma API GraphQL, com schema e query para consultar campos específicos de um jogador.
\end{itemize}

A consulta GraphQL usada no experimento recupera apenas os campos \texttt{nome} e \texttt{gols}. Já o endpoint REST retorna o objeto completo do jogador. Essa diferença representa um caso comum de uso: uma tela ou cliente precisa de poucos campos, mas um endpoint REST retorna mais dados do que o necessário.

\section{Problema foco do experimento}

O problema central é verificar se a adoção de GraphQL, em comparação com REST, traz benefícios mensuráveis em desempenho e volume de dados retornados.

Embora GraphQL seja frequentemente apresentado como uma alternativa mais flexível para evitar \textit{overfetching}, essa vantagem não implica automaticamente menor tempo de resposta. A execução GraphQL inclui parsing da consulta, validação, resolução de campos e serialização da resposta. Portanto, o experimento separa duas dimensões: velocidade e tamanho do payload.

\section{Questões-problema ou Questões-Pesquisa}

\begin{description}[leftmargin=2.5cm]
  \item[RQ1] Respostas às consultas GraphQL são mais rápidas que respostas às consultas REST?
  \item[RQ2] Respostas às consultas GraphQL têm tamanho menor que respostas às consultas REST?
\end{description}

\section{Hipótese(s)}

\subsection{Hipóteses para RQ1}

\begin{itemize}
  \item \textbf{H0\textsubscript{RQ1}:} o tempo médio de resposta do GraphQL é igual ao tempo médio de resposta do REST.
  \item \textbf{H1\textsubscript{RQ1}:} o tempo médio de resposta do GraphQL é menor que o tempo médio de resposta do REST.
\end{itemize}

\subsection{Hipóteses para RQ2}

\begin{itemize}
  \item \textbf{H0\textsubscript{RQ2}:} o tamanho médio da resposta GraphQL é igual ao tamanho médio da resposta REST.
  \item \textbf{H1\textsubscript{RQ2}:} o tamanho médio da resposta GraphQL é menor que o tamanho médio da resposta REST.
\end{itemize}

O nível de significância adotado foi $\alpha = 0{,}05$.

\section{Objetivo}

\subsection{Objetivo principal}

Avaliar quantitativamente se a API GraphQL apresenta respostas mais rápidas e menores que a API REST em um cenário controlado.

\subsection{Objetivos específicos}

\begin{itemize}
  \item implementar APIs REST e GraphQL sobre a mesma base de dados;
  \item executar requisições pareadas para os dois tratamentos;
  \item medir tempo de resposta em milissegundos;
  \item medir tamanho da resposta em bytes;
  \item analisar estatisticamente as diferenças observadas;
  \item gerar tabelas e gráficos para visualização dos resultados;
  \item confrontar os resultados com literatura científica sobre REST e GraphQL.
\end{itemize}

\section{Metodologia}

\subsection{Desenho experimental}

O experimento foi desenhado como um estudo pareado. Em cada iteração, um ID de jogador foi selecionado e usado nos dois tratamentos:

\begin{enumerate}
  \item requisição REST para o endpoint do jogador;
  \item requisição GraphQL para o mesmo jogador, selecionando apenas \texttt{nome} e \texttt{gols}.
\end{enumerate}

Esse desenho reduz a influência da variação entre objetos experimentais, pois REST e GraphQL consultam o mesmo jogador em cada par de medições.

\subsection{Variáveis}

\begin{table}[H]
\centering
\caption{Variáveis do experimento}
\begin{tabular}{lll}
\toprule
Tipo & Variável & Descrição \\
\midrule
Independente & Tecnologia de API & REST ou GraphQL \\
Dependente & \texttt{tempo\_ms} & Tempo de resposta em milissegundos \\
Dependente & \texttt{tamanho\_bytes} & Tamanho do payload em bytes \\
Controle & ID do jogador & Mesmo ID usado no par REST/GraphQL \\
\bottomrule
\end{tabular}
\end{table}

\subsection{Tratamentos}

\begin{itemize}
  \item \textbf{REST:} requisição \texttt{GET} para \texttt{/rest/players/:id}.
  \item \textbf{GraphQL:} requisição \texttt{POST} para \texttt{/graphql} com a query \texttt{player(id) \{ nome gols \}}.
\end{itemize}

\subsection{Objeto experimental}

O objeto experimental é uma requisição de consulta a um jogador existente na base fictícia de estatísticas de futebol.

\section{Passo a passo do experimento}

\begin{enumerate}
  \item Implementação da base de dados em memória.
  \item Implementação da API REST.
  \item Implementação da API GraphQL com Apollo Server.
  \item Implementação do script de validação do ambiente.
  \item Execução de uma coleta de validação com 10 iterações.
  \item Execução da coleta oficial com 1000 iterações por tratamento.
  \item Armazenamento das medições em \texttt{output/results.csv}.
  \item Análise estatística dos dados.
  \item Geração de relatório Markdown.
  \item Geração de gráficos com Matplotlib e Seaborn.
  \item Empacotamento do relatório em projeto LaTeX para importação no Prism.
\end{enumerate}

\section{Decisões}

\begin{itemize}
  \item Usar requisições pareadas para reduzir variação entre objetos.
  \item Usar Docker para melhorar a reprodutibilidade do ambiente.
  \item Manter dependências Python pinadas em \texttt{requirements.txt}.
  \item Medir tamanho da resposta pelo número de bytes retornados.
  \item Usar teste estatístico unicaudal, pois as hipóteses alternativas indicam direção específica: GraphQL menor que REST.
  \item Usar o teste de Mann-Whitney por ser uma alternativa não paramétrica adequada quando não se assume normalidade dos dados.
\end{itemize}

\section{Materiais utilizados}

\begin{table}[H]
\centering
\caption{Materiais e ferramentas}
\begin{tabular}{ll}
\toprule
Material & Uso \\
\midrule
Node.js 20 & Execução da API REST/GraphQL \\
Express & Implementação da API REST \\
Apollo Server & Implementação da API GraphQL \\
Python 3 & Scripts de coleta e análise \\
Requests & Execução das requisições HTTP \\
Pandas & Manipulação dos dados coletados \\
SciPy & Teste estatístico \\
Matplotlib e Seaborn & Geração dos gráficos \\
Docker e Docker Compose & Ambiente controlado de execução \\
Prism & Escrita e colaboração científica em LaTeX \cite{openai2026prism} \\
\bottomrule
\end{tabular}
\end{table}

\section{Métodos utilizados}

O método de coleta mede cada requisição com \texttt{time.perf\_counter}, registra o tempo em milissegundos e calcula o tamanho da resposta com base no conteúdo HTTP retornado. As respostas com erro HTTP, erro GraphQL, timeout ou exceção de rede são classificadas como falhas.

Antes da coleta oficial, o script valida se os endpoints REST e GraphQL respondem corretamente. Se a validação falhar, a coleta oficial é interrompida. Na execução controlada registrada neste relatório, foram obtidos 2000 registros válidos e 0 falhas.

\section{Métricas e suas Unidades}

\begin{table}[H]
\centering
\caption{Métricas avaliadas}
\begin{tabular}{lll}
\toprule
Métrica & Unidade & Relação com a questão de pesquisa \\
\midrule
\texttt{tempo\_ms} & milissegundos (ms) & RQ1: desempenho temporal \\
\texttt{tamanho\_bytes} & bytes & RQ2: volume de dados retornado \\
\bottomrule
\end{tabular}
\end{table}

\section{Visualização dos Resultados}

\subsection{Visualização tabular, lista ou gráfica do que foi colhido no experimento}

A Tabela \ref{tab:descritivas} apresenta as estatísticas descritivas da coleta oficial.

\begin{table}[H]
\centering
\caption{Estatísticas descritivas da coleta oficial}
\label{tab:descritivas}
\begin{tabular}{lrrrrrr}
\toprule
Tecnologia & Métrica & N & Média & Mediana & Desvio padrão & Min--Max \\
\midrule
GraphQL & tempo (ms) & 1000 & 1.7641 & 1.5131 & 1.1198 & 0.8339--19.2746 \\
REST & tempo (ms) & 1000 & 1.2227 & 1.0699 & 0.6051 & 0.6342--9.0806 \\
GraphQL & tamanho (bytes) & 1000 & 54.272 & 54.0 & 2.0000 & 50--59 \\
REST & tamanho (bytes) & 1000 & 326.236 & 326.0 & 2.8111 & 320--333 \\
\bottomrule
\end{tabular}
\end{table}

A Tabela \ref{tab:decisoes} resume as decisões estatísticas para as duas questões de pesquisa.

\begin{table}[H]
\centering
\caption{Resultados estatísticos por questão de pesquisa}
\label{tab:decisoes}
\begin{tabular}{lllll}
\toprule
Questão & Métrica & Diferença das medianas & Valor-p & Decisão \\
\midrule
RQ1 & tempo\_ms & -0.4432 ms & 1.000000 & Não rejeitar H0 \\
RQ2 & tamanho\_bytes & 272 bytes & 0.000000 & Rejeitar H0 \\
\bottomrule
\end{tabular}
\end{table}

Os arquivos completos de dados estão incluídos no projeto:

\begin{itemize}
  \item \texttt{dados/results.csv}: coleta oficial com 2000 registros;
  \item \texttt{dados/validation\_results.csv}: coleta de validação.
\end{itemize}

\section{Gráficos}

\begin{figure}[H]
\centering
\includegraphics[width=0.85\textwidth]{tempo_ms.png}
\caption{Distribuição do tempo de resposta por tecnologia}
\label{fig:tempo}
\end{figure}

\begin{figure}[H]
\centering
\includegraphics[width=0.85\textwidth]{tamanho_bytes.png}
\caption{Distribuição do tamanho da resposta por tecnologia}
\label{fig:tamanho}
\end{figure}

\section{Discussão dos resultados}

\subsection{Confrontar Questões-pesquisa}

\textbf{RQ1.} Os resultados não sustentam a hipótese de que GraphQL foi mais rápido que REST. A mediana do GraphQL foi 1.5131 ms, enquanto a mediana do REST foi 1.0699 ms. A diferença das medianas foi -0.4432 ms quando calculada como REST menos GraphQL, indicando maior tempo para GraphQL. O valor-p foi 1.000000, portanto a hipótese nula não foi rejeitada ao nível de 0.05.

\textbf{RQ2.} Os resultados sustentam a hipótese de que GraphQL teve tamanho de resposta menor que REST. A mediana do GraphQL foi 54 bytes, enquanto a mediana do REST foi 326 bytes. A redução percentual foi 83.44\%, com valor-p igual a 0.000000. Portanto, a hipótese nula de RQ2 foi rejeitada.

\subsection{Insights}

O principal insight é que GraphQL reduziu significativamente o payload porque a consulta pediu apenas dois campos. No entanto, essa redução de bytes não se converteu em menor tempo de resposta no cenário medido. Isso sugere que, para uma API local, pequena e em memória, o custo adicional de processamento GraphQL pode superar a economia de serialização e transferência.

\subsection{Comparações}

REST foi mais rápido no experimento, mas retornou mais dados. GraphQL foi mais econômico em tamanho de resposta, mas teve maior tempo mediano. Assim, a escolha entre REST e GraphQL depende do objetivo principal:

\begin{itemize}
  \item se o objetivo é minimizar payload, GraphQL foi superior neste cenário;
  \item se o objetivo é minimizar latência local para uma consulta simples, REST foi superior neste cenário;
  \item em cenários com múltiplas chamadas REST ou redes mais lentas, a vantagem de GraphQL pode mudar, pois a redução de round trips pode se tornar mais relevante.
\end{itemize}

\subsection{Estatísticas}

A análise estatística usou hipóteses unicaudais. Para RQ1, a alternativa esperava GraphQL menor que REST em tempo. Como os dados apontaram o contrário, o teste não ofereceu evidência para rejeitar H0. Para RQ2, a diferença foi grande e consistente, resultando em rejeição de H0.

\subsection{Confronto com artigos ou trabalhos científicos}

Os resultados de RQ2 corroboram a motivação técnica de GraphQL descrita pela GraphQL Foundation: o cliente seleciona os campos desejados, e a resposta segue o formato da consulta \cite{graphqlqueries2026}. Essa propriedade explica por que o payload GraphQL foi menor neste experimento.

O estudo empírico de Wittern et al. \cite{wittern2019graphqlschemas} também caracteriza GraphQL como uma abordagem em que clientes definem precisamente os dados que desejam recuperar, o que se alinha à redução de resposta observada aqui. Portanto, os resultados de tamanho de payload são consistentes com a literatura.

Por outro lado, os resultados de RQ1 contestam a expectativa de que GraphQL seja necessariamente mais rápido. O trabalho GraphQLify \cite{amareen2026graphqlify} relata ganhos de tempo em cenários que envolvem workflows com múltiplas chamadas REST sequenciais. O presente experimento avaliou uma consulta única e simples, executada localmente contra base em memória. Assim, a divergência é plausível: GraphQL pode reduzir tempo quando evita múltiplos round trips, mas não necessariamente em uma consulta única onde REST já retorna rapidamente.

\section{Conclusão}

\subsection{Tomada de decisão}

Com base nos dados coletados, a adoção de GraphQL seria recomendada quando o objetivo principal for reduzir o tamanho das respostas e evitar retorno de campos desnecessários. Para reduzir tempo de resposta em consultas simples e locais, os resultados deste experimento não justificam substituir REST por GraphQL.

\subsection{Sugestões futuras}

\begin{itemize}
  \item repetir o experimento com banco de dados externo;
  \item testar consultas que exigem múltiplos endpoints REST;
  \item medir também uso de CPU e memória;
  \item avaliar impacto de cache;
  \item executar os tratamentos em ordem aleatorizada;
  \item testar com maior latência de rede;
  \item comparar consultas GraphQL com diferentes quantidades de campos.
\end{itemize}

\subsection{Resultado conclusivo sucinto}

Neste experimento controlado, GraphQL reduziu o tamanho das respostas em aproximadamente 83.44\%, mas não reduziu o tempo de resposta. Portanto, a principal vantagem observada foi economia de payload, não desempenho temporal.

\section{Ameaças à validade}

\begin{itemize}
  \item \textbf{Validade interna:} ordem das requisições, aquecimento do runtime, cache e carga da máquina podem influenciar tempos.
  \item \textbf{Validade externa:} a base é fictícia e em memória; os resultados podem não generalizar para APIs reais com banco externo.
  \item \textbf{Validade de construção:} tempo e tamanho de resposta não capturam todos os custos de adoção de GraphQL.
  \item \textbf{Validade de conclusão:} embora a amostra tenha 1000 registros por tratamento, novas execuções podem variar nos tempos absolutos.
\end{itemize}

\bibliographystyle{plain}
\bibliography{referencias}

\end{document}
